\section{Define which error wins}
Several defects can coexist. A header might have both the wrong magic and an
unsupported version. A declared length of 65535 might also exceed the available
input. The interface defines an order so that independently written parsers return
the same status. Apply the following decisions from top to bottom.

\Needspace{14\baselineskip}
\begin{longtable}{rlp{91mm}}
  \toprule Status & Name         & Condition, after earlier checks pass                 \\\midrule
  \(-1\)          & Need more    & Fewer than eight available bytes                     \\
  \(-2\)          & Bad magic    & Either identifier byte differs                       \\
  \(-3\)          & Bad version  & Version byte is not one                              \\
  \(-4\)          & Bad flags    & Any reserved flag bit is set                         \\
  \(-5\)          & Bad length   & Declared payload length exceeds 64                   \\
  \(-1\)          & Need more    & Fewer than \(9+L\) bytes are available               \\
  \(-6\)          & Bad checksum & Sum of the first \(9+L\) bytes is nonzero modulo 256 \\
  \bottomrule
\end{longtable}
For seven available bytes, return \texttt{-1} even if the first byte is wrong.
The format checks begin only once a full fixed header is available. For eight
bytes with correct earlier fields and a declared length of 65, return \texttt{-5}.
The semantic length limit takes priority over missing body bytes. For a valid
length of three with eleven available bytes, return \texttt{-1}; the twelfth byte
is still required as the trailer.

This ordering constrains returned behavior, not every internal load. An alternative
parser can read all eight header bytes after confirming their availability and then
validate their fields. It can also use a wider aligned load wholly within the
supplied interval. Grading does not require the demonstration's exact read order or
number of loads. It does require the specified error and safe complete read widths.
